\section{Постановка и программа исследования}\label{sec:program}

Во всём приложении \emph{выпрямлением} называется автоморфизм всего
аффинного пространства, переводящий заданную параметризованную кривую
в стандартную прямую. Выбор поля и требование полиномиальности обратного
отображения входят в постановку. Бирациональное выпрямление является
более слабым условием.

\subsection{Координатный подъём параметра}

Пусть $k$ --- поле, $R=k[x,y,z]$ и
\[
 \gamma(t)=(p(t),q(t),r(t)),\qquad p,q,r\in k[t].
\]
Параметризация задаёт замкнутое вложение $\A^1_k\hookrightarrow\A^3_k$
тогда и только тогда, когда
\begin{equation}\label{eq:closed-embedding}
 k[p,q,r]=k[t].
\end{equation}
Это эквивалентно сюръективности гомоморфизма ограничения
$\gamma^*:R\to k[t]$. В частности, существует полином $H_0\in R$
с $H_0\gamma=t$. Однако такой полином ещё не обязан быть
\emph{координатой}, то есть компонентой полиномиального автоморфизма.

\begin{lemma}[Координатный подъём]\label{lem:coordinate-lift}
Вложение $\gamma$ допускает выпрямитель $\Phi\in\Aut_k(\A^3)$,
$\Phi\gamma(t)=(t,0,0)$, тогда и только тогда, когда существует
координата $H\in R$, удовлетворяющая $H\gamma=t$.
\end{lemma}
\begin{proof}
Необходимость даёт первая компонента $\Phi$. Обратно, дополним $H$
до автоморфизма $(H,U,V)$. На кривой его значение равно
$(t,u(t),v(t))$. Применим затем автоморфизм
\[
 (T,Y,Z)\longmapsto(T,Y-u(T),Z-v(T)).
\]
Полученная композиция выпрямляет $\gamma$.
\end{proof}
Критерий известен; он явно используется в \cite[лемма~1]{Lewis2024}.
Его конструктивное содержание состоит в выборе координаты среди всех
подъёмов $H_0+\ker\gamma^*$.

\subsection{Что даёт двумерный случай}

В характеристике нуль теорема Абъянкара--Моха--Судзуки утверждает,
что всякое замкнутое вложение $\A^1_k\hookrightarrow\A^2_k$
выпрямляется \cite{AM1975,Suzuki1974}. Для непараметризованных образов
это означает транзитивность действия $G_2(k)=\Aut_k(\A^2)$.
Положив $L_2=\{y=0\}$, получаем описание множества таких образов:
\[
 \{\text{образы замкнутых вложений }\A^1_k\hookrightarrow\A^2_k\}
 \simeq G_2(k)/\operatorname{Stab}(L_2).
\]
Здесь используются левые смежные классы $g\operatorname{Stab}(L_2)$, причём
\begin{equation}\label{eq:plane-stabilizer}
 \operatorname{Stab}(L_2)=
 \{(x,y)\mapsto(ax+P(y),by):a,b\in k^*,\ P\in k[y]\}.
\end{equation}
Действительно, сохранение идеала $(y)$ даёт образ $y$, равный $by$.
После масштабирования остаётся автоморфизм $k[y][x]$, поэтому образ $x$
имеет вид $ax+P(y)$ с $a\in k^*$. Обратное включение непосредственно.

Следует различать три классических утверждения. Теорема
Абъянкара--Моха--Судзуки описывает орбиту вложенных прямых. Теорема
Юнга--ван дер Кулка описывает $G_2(k)$ через аффинные и треугольные
автоморфизмы и справедлива над любым полем
\cite[теорема~1]{EV2001}. Теорема Нётера--Кастельнуово утверждает
порождение $\Cr_2(\C)$ проективными автоморфизмами и стандартной
квадратичной инволюцией; см. обсуждение в \cite[\S1]{BH2015}.
Описание орбиты прямой само по себе не является описанием всей группы
Кремоны. В положительной характеристике общая плоская теорема о
выпрямлении неверна, поэтому перенос её заключения на произвольное
поле без оговорок недопустим.

\subsection{Бирациональное выпрямление и точная групповая постановка}

Выбор стандартной аффинной карты даёт включение
\[
 G=\Aut_k(\A^3)\subset B=\Cr_3(k)=\operatorname{Bir}_k(\mathbb P^3).
\]
Для $b\in B$ имеется точный критерий:
\begin{equation}\label{eq:ring-preservation}
 b\in G\quad\Longleftrightarrow\quad b^*R=R.
\end{equation}
Действительно, равенство означает полиномиальность всех компонент $b$
и $b^{-1}$. В терминах порядков вдоль неприводимых полиномов $P\in R$
это условие записывается как
\[
 v_P(b_i)\geq0,\qquad v_P((b^{-1})_i)\geq0
 \quad\text{для всех }P\text{ и }i.
\]
Используется факториальность $R$. Полюса на гиперплоскости бесконечности
при этом допустимы. Проверка одного только якобиана не заменяет
условия~\eqref{eq:ring-preservation}.

\begin{proposition}\label{prop:birational-any-field}
Над любым полем $k$ всякое вложение~\eqref{eq:closed-embedding}
имеет бирациональный выпрямитель над $k$. Его и обратное можно выбрать
регулярными около общих точек соответствующих кривых. Рациональный
якобиан выпрямителя можно сделать равным $1$.
\end{proposition}
\begin{proof}
Дифференцирование равенства $H_0(p,q,r)=t$ даёт
\[
 (H_0)_x(p,q,r)p'+(H_0)_y(p,q,r)q'+(H_0)_z(p,q,r)r'=1.
\]
Поэтому после циклической перестановки координат можно считать
$p'\ne0$. Расширение $k(t)/k(p)$ конечно и сепарабельно, и
$k(t)=k(p)(q,r)$. По теореме о примитивном элементе можно выбрать
$\lambda(p)\in k[p]$ так, чтобы
$k(p,q+\lambda(p)r)=k(t)$. В самом деле, плохих коэффициентов
$\lambda\in k(p)$ лишь конечное число, а $k[p]$ бесконечно даже
при конечном $k$. Это следует также непосредственно из сравнения
различных вложений сепарабельного расширения: каждое неравное
сопряжение исключает не более одного значения $\lambda$.

Применим сдвиг $(x,y,z)\mapsto(x,y+\lambda(x)z,z)$ и снова обозначим
вторую компоненту кривой через $q$. Существует
$\tau\in k(x,y)$ с $\tau(p,q)=t$; знаменатель выбранного выражения
не равен нулю тождественно на кривой. Положим
\begin{equation}\label{eq:birational-construction}
 s=z-r(\tau(x,y))+\tau(x,y),\qquad
 b=(s,x-p(s),y-q(s)).
\end{equation}
Тогда $b\gamma=(t,0,0)$. Для входа $(T,U,V)$ обратное задаётся формулами
\begin{equation}\label{eq:birational-inverse}
 \begin{split}
 X&=U+p(T),\qquad Y=V+q(T),\\
 b^{-1}(T,U,V)&=(X,Y,T+r(\tau(X,Y))-\tau(X,Y)).
 \end{split}
\end{equation}
Обе композиции тождественны непосредственной подстановкой. На прямой
$(T,0,0)$ знаменатель становится тем же ненулевым элементом $k(T)$,
поэтому обратное также определено в общей точке. Все использованные
сдвиги имеют якобиан $1$, как и циклические перестановки.
\end{proof}

Над алгебраически замкнутым полем бирациональная эквивалентность
проективного замыкания нашей кривой и прямой является частным случаем
теоремы Меллы--Поластри: два бирациональных вложения $d$-мерного
многообразия в $\mathbb P^n$ кремоново эквивалентны при $n\geq d+2$
\cite[теорема~1]{MP2009}. Приведённое элементарное доказательство
указывает карту над исходным полем. Ни одно из этих утверждений
не гарантирует регулярность во всех точках аффинного пространства.

Для Шастри можно явно взять
\[
 \tau(x,y)=\frac{x(y+1)}{x^2-y-3},\qquad
 s=z-\tau^5+11\tau,\qquad
 b=(s,x-s^3+3s,y-s^4+4s^2).
\]
Равенство $p(q+1)=t(p^2-q-3)$ при $p=t^3-3t$, $q=t^4-4t^2$
проверяет восстановление параметра. Это бирациональный выпрямитель
над $\Q$, допускающий полюса; обе его композиции с явным обратным
проверяются в новом вычислительном пакете. Его существование совместимо
с невозможностью полиномиального выпрямления над $\Q$, обсуждаемой ниже.

Пусть $L=\{y=z=0\}$. Обозначим через $S_L$ подгруппу $B$,
состоящую из бирациональных карт, задающих изоморфизмы между некоторыми
окрестностями общей точки $L$ и действующих на $L$ тождественно.
Таким образом, условие определённости предъявляется и к обратной карте.

\begin{proposition}[Задача пересечения смежного класса]\label{prop:coset}
Если $b$ --- бирациональный выпрямитель из
предложения~\ref{prop:birational-any-field}, то
\begin{equation}\label{eq:coset}
 \gamma\text{ полиномиально выпрямляется}
 \quad\Longleftrightarrow\quad S_Lb\cap G\ne\varnothing.
\end{equation}
\end{proposition}
\begin{proof}
Для полиномиального выпрямителя $\Phi$ карта $h=\Phi b^{-1}$
принадлежит $S_L$, и $\Phi=hb$. Обратно, если $hb\in G$ при
$h\in S_L$, то $(hb)\gamma(t)=(t,0,0)$ как рациональное равенство.
Поскольку обе стороны полиномиальны, равенство справедливо всюду.
\end{proof}
Таким образом, исходная задача переформулируется как поиск поправки
$h$, устраняющей аффинные полюса $b$ и его обратного при сохранении
параметризованной прямой. Формула~\eqref{eq:coset} сама по себе не
является алгоритмом и не доказывает существование нужного пересечения.

\subsection{Относительная подгруппа как рабочий механизм}

Стабилизатор рациональной функции $x$ и его полиномиальная часть равны
\[
 J_x=\{b\in B:b^*x=x\}\simeq\Cr_2(k(x)),\qquad
 G_x=G\cap J_x=\Aut_{k[x]}k[x,y,z].
\]
Элемент $G_x$ становится ручным автоморфизмом двух переменных над
$k(x)$, но его факторизация над этим полем может иметь знаменатели.
Их устранение над $k[x]$ --- содержательная часть конструкции.
Этот подход имеет прямых предшественников в теории переноса координат
и остаточных координат \cite{EV2001,Lewis2024}.

Если $A=(x,H,W)\in G_x$, $H\gamma=t$ и $W\gamma=w(t)$, то
\begin{equation}\label{eq:relative-rectifier}
 \Phi=(H,x-p(H),W-w(H))
\end{equation}
выпрямляет $\gamma$. В самой группе $G_x$ естественная целевая
параметризация имеет вид $(p(t),t,0)$: фиксирующие $x$ карты не меняют
первую компоненту $p(t)$. Последующий внешний ручной автоморфизм
переводит её в $(t,0,0)$.

\begin{remark}\label{rem:relative-not-necessary}
Для заранее фиксированного $x$ существование относительной координаты
$H$ с $H\gamma=t$ является достаточным условием выпрямляемости.
Его необходимость для произвольного выпрямляемого вложения здесь не
доказана. Разрешение произвольной предварительной полиномиальной
замены без дополнительных ограничений сделало бы утверждение о
полноте круговым: такой заменой мог бы оказаться искомый выпрямитель.
\end{remark}

\subsection{Поле определения и препятствие к спуску}

Для кривой Шастри
\[
 \gamma_T=(t^3-3t,t^4-4t^2,t^5-10t)
\]
вещественный образ задаёт трилистник \cite[\S3, пример~1]{Shastri1992}.
Полиномиальный автоморфизм $\mathbb R^3$ является собственным
диффеоморфизмом и продолжается до гомеоморфизма одноточечной
компактификации $S^3$. Он не переводит нетривиальный узел в тривиальный.
Поэтому выпрямление над $\mathbb R$, а также над $\Q$, невозможно.
Тот же аргумент применим к любому числовому полю с вещественным
вложением: достаточно применить его к коэффициентам предполагаемой карты.

В сочетании с явным выпрямителем над $K=\Q(\sqrt{-3})$ это даёт
\[
 \min\{[F:\Q]: F\text{ --- числовое поле, над которым }
 \gamma_T\text{ выпрямляется}\}=2.
\]
Подходящее числовое поле должно быть вполне мнимым. Единственность
$K$ среди квадратичных полей не установлена; отсутствие в поле корня
$a^2-a+1$ исключает непосредственную специализацию нашей формулы,
но не другую конструкцию.

Полиномиальный спуск допускает точную групповую запись. Пусть
$\alpha\in\Aut_K(\A^3)$ выпрямляет $\gamma_T$,
$\Gamma=\operatorname{Gal}(K/\Q)$ и $H_L(K)$ --- поточечный
\emph{полиномиальный} стабилизатор $L$. Тогда
\[
 c_\sigma=\alpha\,\sigma(\alpha)^{-1}\in H_L(K),\qquad
 c_{\sigma\tau}=c_\sigma\,\sigma(c_\tau).
\]
Выпрямитель над $\Q$ существует тогда и только тогда, когда найдётся
$h\in H_L(K)$ с $c_\sigma=h^{-1}\sigma(h)$ для всех $\sigma$.
Действительно, это равносильно определённости $h\alpha$ и его
полиномиального обратного над $\Q$. Для Шастри такой $h$ не существует.
После замены полиномиального стабилизатора на $S_L$ препятствие исчезает:
при бирациональном выпрямителе $b$ над $\Q$ подходит $h=b\alpha^{-1}$.
Следовательно, переход к полной группе Кремоны теряет в том числе
информацию о полиномиальном спуске. Здесь используется лишь
элементарное равенство кокциклов; классификация некоммутативного
множества $H^1(\Gamma,H_L(K))$ не заявляется.

\subsection{Что остаётся программой, а не результатом}

\begin{enumerate}
\item \textbf{Полнота переносов.} Описать класс параметризаций,
для которого выбор центров модификаций и исключительных дивизоров
всегда завершается относительной координатой. Конечные вычисления
для отдельных индексов не доказывают такую полноту.
\item \textbf{Контролируемые изменения проекции.} Выяснить, какие
заранее заданные аффинные или ограниченные полиномиальные замены
позволяют расширить применимость $G_x$ без круговой постановки.
\item \textbf{Ручная орбита самой кривой.} Дикость построенного
выпрямителя не исключает другого ручного выпрямителя. Требуется
исследовать все координаты в $H+\ker\gamma^*$, а не только выбранную $H$.
\item \textbf{Арифметика.} Определить другие поля выпрямления,
например проверить существование выпрямителя Шастри над $\Q(i)$.
Отсутствие корня в конкретной формуле и препятствие к существованию
любого автоморфизма --- разные утверждения.
\end{enumerate}
Эти цели не требуют предварительного полного описания $\Cr_3(\C)$.
Они исследуют её определённые подгруппы, смежные классы и условия
регулярности на фиксированной аффинной модели.
